% =============================================================================
%  Lab 7 - Next Instruction Address Generation I
% =============================================================================

\chapter{Lab 7 --- Next Instruction Address Generation I}

\section{Objective}
% Two or three sentences, your own words.

\section{Specification}

\begin{table}[htbp]
\centering
\caption{Transition table (Type T) for \texttt{NIAG\_Reg}, the program counter.}
\label{tab:lab07-truth}
\begin{tabular}{cccl||l}
\toprule
\# & \texttt{Rst} & \texttt{Strb} & \texttt{Clk} & $PC^{+}$ \\
\midrule
0 & 1 & --- & --- & \texttt{RESET\_VEC} (asynchronous) \\
1 & 0 & --- & not $\uparrow$ & $PC$ --- hold \\
2 & 0 & 0 & $\uparrow$ & $PC$ --- edge, not loading (a stall) \\
3 & 0 & 1 & $\uparrow$ & \texttt{NextPC} \\
\bottomrule
\end{tabular}
\end{table}

\begin{align*}
PC^{+} &= \overline{\texttt{Rst}} \cdot \left( \texttt{Strb}\cdot\texttt{NextPC} + \overline{\texttt{Strb}}\cdot PC \right) \\
\texttt{NSA} &= PC + 4 \\
\texttt{Rel} &= \texttt{NSA} + \texttt{Xtnd32}
\end{align*}

% ONE SENTENCE, required: compare the PC equation with your Lab 2 register
% equation. What actually changed?

\section{Design}
% Answer all three:
%   (a) Why is the branch base PC+4 and not PC? Answer as a trade-off, not a
%       rule - what would the assembler have to do if the base were PC?
%   (b) Why is the branch offset shifted left by two? How much branch range
%       does that buy, and what makes it safe?
%   (c) Did you write NIAG_Reg out, or instantiate your Lab 2 RegN? If you
%       reused RegN, how did you handle the reset VALUE (RESET_VEC) rather
%       than reset-to-zero?

\section{Implementation}

\vhdlfile{code/lab07/NIAG_NSA.vhd}{\texttt{NIAG\_NSA} --- next sequential address}

\vhdlfile{code/lab07/NIAG_Rel.vhd}{\texttt{NIAG\_Rel} --- PC-relative branch target}

\vhdlfile{code/lab07/NIAG_Reg.vhd}{\texttt{NIAG\_Reg} --- the program counter}

\section{Verification}

\subsection*{Hand-computed arithmetic}
% REQUIRED. Do this BEFORE simulating, for at least one FORWARD and one
% BACKWARD branch. Show all three steps each time, then put the simulated
% value beside it.

\begin{table}[htbp]
\centering
\caption{Branch targets: computed by hand, then by the hardware.}
\label{tab:lab07-hand}
\begin{tabular}{llllll}
\toprule
$PC$ & immediate & \texttt{NSA} & offset & target (by hand) & target (simulated) \\
\midrule
\texttt{0x0000000C} & \texttt{0x0002} & & & & \\
 & & & & & \\
\bottomrule
\end{tabular}
\end{table}

\subsection*{Test vectors and coverage}
% The coverage argument here is DIRECTIONS AND BOUNDARIES:
%   forward, backward, zero offset, largest positive, most negative.
% Say explicitly which bug an all-positive test set would fail to detect.

\section{Results}
% Evidence: E1, E2, E4, E7.

\figplaceholder{the waveform showing a forward branch target and a backward
branch target, both annotated with their computed values}

\figplaceholder{the RTL Viewer, showing the PC register and the two adders}

% E4: register count. Expect 32 - one per PC bit. NIAG_NSA and NIAG_Rel are
%     pure adders and should contain no storage at all. Say so.

\section{Discussion}
% REQUIRED. One bug, how you found it, what it cost.
%
% Good honest topic if you need one: did your first branch target come out
% wrong? Was it off by a FACTOR (forgot the shift) or off by a CONSTANT (used
% PC as the base)? How did you tell which - and would you have noticed the
% second one if the first had not happened?
